\documentclass[11pt]{article}
\usepackage[margin=1in]{geometry}
\usepackage[T1]{fontenc}
\usepackage[utf8]{inputenc}
\usepackage{lmodern}
\usepackage{microtype}
\usepackage{booktabs}
\usepackage{tabularx}
\usepackage{array}
\usepackage[table]{xcolor}
\usepackage{hyperref}
\usepackage{verbatim}
\definecolor{TitleBlue}{HTML}{0E3A5D}
\definecolor{MetaGray}{HTML}{58636E}
\definecolor{Aligned}{HTML}{2E7D32}
\definecolor{Partial}{HTML}{F57F17}
\definecolor{Contradicted}{HTML}{C62828}
\definecolor{Unclear}{HTML}{616161}
\definecolor{NewFind}{HTML}{1565C0}
\hypersetup{
  colorlinks=true,
  linkcolor=TitleBlue,
  urlcolor=TitleBlue,
  citecolor=TitleBlue
}
\setlength{\parindent}{0pt}
\setlength{\parskip}{0.45em}
\renewcommand{\arraystretch}{1.18}
\newcolumntype{Y}{>{\raggedright\arraybackslash}X}
\begin{document}
\begin{center}
{\LARGE \textbf{\textcolor{TitleBlue}{Integrated Experiment Report}}}\\[0.25em]
{\Large \textbf{Mug Bargain Cheap Talk}}\\[0.45em]
{\color{MetaGray} Generated on 2026-02-23 \quad$\cdot$\quad N=360 sessions \quad$\cdot$\quad 12 conditions \quad$\cdot$\quad gpt-4o}
\end{center}
\vspace{0.4em}
\hrule
\vspace{0.7em}
\textbf{Experiment ID:} \texttt{mug\_bargain\_cheap\_talk}\\
\textbf{Experiment Directory:} \texttt{experiments/mug\_bargain\_cheap\_talk}\\
\textbf{Model:} gpt-4o \quad \textbf{Design:} $3 \times 2 \times 2$ full factorial \quad \textbf{N per cell:} 30
\vspace{0.6em}

%% ============================================================
\section{Research Question}
%% ============================================================

How does pre-play cheap talk affect bilateral bargaining efficiency over a mug, and how does the effect depend on (a) the degree of preference misalignment between buyer and seller and (b) the information structure?

This experiment is grounded in two foundational papers:
\begin{itemize}
\item \textbf{Crawford \& Sobel (1982)} --- strategic information transmission via partition equilibria; the bias parameter $b$ governs equilibrium informativeness.
\item \textbf{Farrell \& Gibbons (1989)} --- cheap talk in extensive-form games; pooling equilibria always exist, separating equilibria require neologism-proofness.
\end{itemize}

%% ============================================================
\section{Prior Paper Findings and Experiment Settings}
%% ============================================================

Twelve design choices were extracted from the two source papers and recorded in \texttt{evidence\_trace.csv}. The key theoretical predictions are:

\begin{table}[htbp]
\centering
\small
\begin{tabularx}{\textwidth}{@{}l l Y@{}}
\toprule
\textbf{ID} & \textbf{Source} & \textbf{Prior claim} \\
\midrule
CS82\_partition & CS82 & Cheap talk transmits partial information via partition equilibria; more communication rounds $\to$ finer partitions $\to$ better outcomes \\
CS82\_bias & CS82 & Smaller $b$ (narrow gap) $\to$ finer partitions $\to$ more informative talk \\
CS82\_aligned & CS82 & Talk more effective when interests more aligned \\
CS82\_sender & CS82 & One-sided info maps to pure sender-receiver game; talk should help more \\
CS82\_uninformative & CS82 & As $b \to \infty$, only uninformative equilibrium remains \\
FG89\_costless & FG89 & Costless non-binding communication should weakly improve outcomes \\
FG89\_pooling & FG89 & Pooling equilibrium always exists; talk may be uninformative \\
FG89\_separating & FG89 & Separating equilibrium: receiver learns sender type \\
FG89\_neologism & FG89 & Neologism-proofness favors separating when it exists \\
\bottomrule
\end{tabularx}
\caption{Key prior claims from Crawford \& Sobel (1982) and Farrell \& Gibbons (1989).}
\label{tab:prior-claims}
\end{table}

\paragraph{Referenced papers (manifest):}
\begin{itemize}
\item \texttt{cheap\_talk}: Crawford \& Sobel (1982), \textit{Econometrica}
\item \texttt{strategic}: Farrell \& Gibbons (1989), \textit{Bell Journal of Economics}
\end{itemize}

%% ============================================================
\section{New Experiment Design and Innovations}
%% ============================================================

\paragraph{Factorial design.} $3 \times 2 \times 2 = 12$ conditions crossing:
\begin{itemize}
\item \textbf{Communication:} none, cheap\_talk (1 round), extended\_talk (3 rounds)
\item \textbf{Value gap:} narrow (overlapping WTP/WTA), wide (non-overlapping $\sim$50\%)
\item \textbf{Info asymmetry:} two\_sided\_private, one\_sided\_buyer\_known
\end{itemize}

\paragraph{Innovations vs.\ prior work.}
\begin{enumerate}
\item \textbf{LLM agents as subjects} --- eliminates experimenter demand effects, enables 360-session production run
\item \textbf{Three communication levels} --- distinguishes null talk from extended talk (prior work tests binary)
\item \textbf{Holdout validation} --- C10 and C12 withheld from model fitting to test out-of-sample prediction
\item \textbf{Concordia framework} --- structured action parsing, payoff verification, guardrail enforcement
\end{enumerate}

\begin{table}[htbp]
\centering
\small
\begin{tabularx}{\textwidth}{@{}l l l l c r@{}}
\toprule
\textbf{Condition} & \textbf{Communication} & \textbf{Value gap} & \textbf{Info} & \textbf{Holdout} & \textbf{N} \\
\midrule
C01 & none & narrow & two\_sided & & 30 \\
C02 & none & narrow & one\_sided & & 30 \\
C03 & none & wide & two\_sided & & 30 \\
C04 & none & wide & one\_sided & & 30 \\
C05 & cheap\_talk & narrow & two\_sided & & 30 \\
C06 & cheap\_talk & narrow & one\_sided & & 30 \\
C07 & cheap\_talk & wide & two\_sided & & 30 \\
C08 & cheap\_talk & wide & one\_sided & & 30 \\
C09 & extended\_talk & narrow & two\_sided & & 30 \\
C10 & extended\_talk & narrow & one\_sided & $\checkmark$ & 30 \\
C11 & extended\_talk & wide & two\_sided & & 30 \\
C12 & extended\_talk & wide & one\_sided & $\checkmark$ & 30 \\
\bottomrule
\end{tabularx}
\caption{Full $3 \times 2 \times 2$ condition matrix (N=360 total).}
\label{tab:conditions}
\end{table}

%% ============================================================
\section{Run Results Analysis}
%% ============================================================

All numeric claims trace to \texttt{analysis/summary\_by\_condition.csv} and \texttt{analysis/model\_results.csv}.

\subsection{Condition-level deal rates}

\begin{table}[htbp]
\centering
\small
\begin{tabularx}{\textwidth}{@{}l r r r r r@{}}
\toprule
\textbf{Condition} & \textbf{N} & \textbf{Deal \%} & \textbf{95\% CI} & \textbf{Avg price} & \textbf{Avg surplus} \\
\midrule
C01 none/narrow/2s & 30 & 66.7 & [48.8, 80.8] & \$7.95 & 0.64 \\
C02 none/narrow/1s & 30 & 90.0 & [74.4, 96.5] & \$7.11 & 0.73 \\
C03 none/wide/2s & 30 & 30.0 & [16.7, 47.9] & \$8.67 & 0.52 \\
C04 none/wide/1s & 30 & 66.7 & [48.8, 80.8] & \$7.25 & 0.63 \\
C05 talk/narrow/2s & 30 & 63.3 & [45.5, 78.1] & \$8.63 & 0.57 \\
C06 talk/narrow/1s & 30 & 96.7 & [83.3, 99.4] & \$7.55 & 0.65 \\
C07 talk/wide/2s & 30 & 50.0 & [33.2, 66.8] & \$8.33 & 0.55 \\
C08 talk/wide/1s & 30 & 43.3 & [27.4, 60.8] & \$7.38 & 0.68 \\
C09 ext/narrow/2s & 30 & 43.3 & [27.4, 60.8] & \$7.85 & 0.55 \\
C10 ext/narrow/1s & 30 & 100.0 & [88.7, 100] & \$7.37 & 0.70 \\
C11 ext/wide/2s & 30 & 36.7 & [21.9, 54.5] & \$7.91 & 0.56 \\
C12 ext/wide/1s & 30 & 46.7 & [30.2, 63.9] & \$7.45 & 0.61 \\
\bottomrule
\end{tabularx}
\caption{Condition-level results (N=30 per cell). CI = Wilson score 95\% interval.}
\label{tab:results-condition}
\end{table}

\subsection{Main effects (marginal)}

\begin{table}[htbp]
\centering
\small
\begin{tabular}{@{}l r l r@{}}
\toprule
\textbf{Factor} & \textbf{Effect (pp)} & \textbf{Direction} & \textbf{$p$-value} \\
\midrule
Value gap (narrow $-$ wide) & $+31.1$ & narrow $\gg$ wide & $< 0.0001$ \\
Info (one-sided $-$ two-sided) & $+25.8$ & one-sided $\gg$ two-sided & $< 0.0001$ \\
Communication (any talk $-$ none) & $-7.8$ & talk $<$ none & $0.190$ \\
\quad cheap\_talk $-$ none & $+0.0$ & null & $1.000$ \\
\quad extended\_talk $-$ none & $-23.3$ & extended $\ll$ none & $\mathbf{0.004}$ \\
\bottomrule
\end{tabular}
\caption{Marginal main effects on deal rate. Bold = significant at $\alpha = 0.05$.}
\label{tab:main-effects}
\end{table}

\subsection{Hypothesis tests}

\begin{table}[htbp]
\centering
\small
\begin{tabularx}{\textwidth}{@{}l r r l Y@{}}
\toprule
\textbf{Hyp} & \textbf{Effect} & \textbf{$p$} & \textbf{Verdict} & \textbf{Note} \\
\midrule
H1: Talk $\to$ deals & $-7.8$pp & 0.190 & REFUTED & Pooled talk vs none; n.s.\ after Holm ($p_{\text{adj}}=0.380$) \\
H1a: 1-round vs none & $+0.0$pp & 1.000 & NULL & Exactly zero effect \\
H1b: Extended vs none & $-23.3$pp & \textbf{0.004} & REFUTED & Extended talk \emph{harmful}; significant \\
H2: Talk $\times$ gap & $-5.6$pp & 0.480 & REFUTED & No interaction; Holm $p=0.480$ \\
H3: Surplus capture & $-10.4$pp & 0.074 & UNCLEAR & Marginal; reversed direction \\
H4: Info $\times$ talk & $-8.3$pp & 0.393 & REFUTED & One-sided info does not amplify talk \\
H5: Ext $>$ cheap vol. & 6.0 vs 2.0 msgs & $<0.001$ & CONFIRMED & More messages but worse outcomes \\
H6: Price fairness & 0.32 vs 0.21 & 0.399 & UNCLEAR & No convergence toward 50-50 \\
\bottomrule
\end{tabularx}
\caption{Hypothesis test results. Holm-Bonferroni applied to H1--H3.}
\label{tab:hypotheses}
\end{table}

\subsection{Effect size ranking}

The three factors ranked by absolute effect on deal rate:
\begin{enumerate}
\item \textbf{Value gap:} $+31.1$pp ($p < 0.0001$) --- structural feasibility dominates
\item \textbf{Information structure:} $+25.8$pp ($p < 0.0001$) --- one-sided info creates common ground
\item \textbf{Extended talk:} $-23.3$pp ($p = 0.004$) --- communication backfires
\item \textbf{Cheap talk:} $+0.0$pp ($p = 1.000$) --- exactly null
\end{enumerate}

\subsection{Holdout validation}

Two conditions (C10, C12) were excluded from the logistic model and predicted out-of-sample:

\begin{table}[htbp]
\centering
\begin{tabular}{@{}l r r r l@{}}
\toprule
\textbf{Condition} & \textbf{Predicted} & \textbf{Observed} & \textbf{Residual} & \textbf{Status} \\
\midrule
C10 (ext/narrow/1s) & 0.842 & 1.000 & $+0.158$ & Miscalibrated \\
C12 (ext/wide/1s) & 0.610 & 0.467 & $-0.144$ & Calibrated \\
\bottomrule
\end{tabular}
\caption{Holdout predictions vs.\ observed deal rates.}
\label{tab:holdout}
\end{table}

The model underpredicts C10 (ceiling effect at 100\%) and slightly overpredicts C12. Average absolute residual = 0.151.

%% ============================================================
\section{Alignment with Prior Findings and New Discoveries}
%% ============================================================

Each prior claim from Crawford \& Sobel (1982) and Farrell \& Gibbons (1989) is evaluated against our observed results. Full mapping in \texttt{analysis/prior\_alignment.csv}.

\begin{table}[htbp]
\centering
\small
\begin{tabularx}{\textwidth}{@{}l l l Y@{}}
\toprule
\textbf{Prior claim} & \textbf{Source} & \textbf{Status} & \textbf{Evidence} \\
\midrule
Partition signaling & CS82 & \textcolor{Contradicted}{\textbf{CONTRADICTED}} & 0.0pp effect; agents produce rapport not value signals \\
Bias $\to$ finer partitions & CS82 & \textcolor{Contradicted}{\textbf{CONTRADICTED}} & Narrow $-10.6$pp vs wide $-5.0$pp; interaction n.s. \\
Aligned interests $\to$ better talk & CS82 & \textcolor{Contradicted}{\textbf{CONTRADICTED}} & Talk effect slightly more negative when aligned \\
One-sided $\to$ sender-receiver & CS82 & \textcolor{Contradicted}{\textbf{CONTRADICTED}} & $-8.3$pp interaction; one-sided does not amplify talk \\
$b \to \infty$ uninformative & CS82 & \textcolor{Aligned}{\textbf{ALIGNED}} & Wide gap: 45.6\% deal rate; $\sim$50\% infeasible \\
Extended $\to$ finer partition & CS82 & \textcolor{Contradicted}{\textbf{CONTRADICTED}} & Extended: 40.0\% vs none: 63.3\% ($p=0.004$); harmful \\
Costless talk weakly helps & FG89 & \textcolor{Contradicted}{\textbf{CONTRADICTED}} & 0.0pp for 1-round; $-23.3$pp for extended ($p=0.004$) \\
Pooling always exists & FG89 & \textcolor{Aligned}{\textbf{ALIGNED}} & All communication is generic rapport; no value signals \\
Separating equilibrium & FG89 & \textcolor{Contradicted}{\textbf{CONTRADICTED}} & No evidence of value signal transmission \\
Neologism-proofness & FG89 & \textcolor{Contradicted}{\textbf{CONTRADICTED}} & No separating behavior despite favorable conditions \\
Mutual discipline & FG89 & \textcolor{Unclear}{\textbf{UNCLEAR}} & Not testable in bilateral setting \\
\bottomrule
\end{tabularx}
\caption{Prior alignment summary: 2 ALIGNED, 7 CONTRADICTED, 1 UNCLEAR, 1 PARTIAL (not shown: CS82 info similarity).}
\label{tab:prior-alignment}
\end{table}

\subsection{Summary verdict}

The dominant pattern is \textbf{pooling}: LLM agents universally play the pooling equilibrium predicted by FG89 as always existing. They \emph{never} partition-signal or separate, contradicting CS82's core prediction that cheap talk transmits partial information. The practical consequence is that communication is either null (1-round) or actively harmful (extended), likely through anchoring mechanisms where early price mentions create rigid expectations.

\subsection{New discoveries}

\begin{enumerate}
\item \textbf{Information structure dominance} ($+25.8$pp, $p < 0.0001$): One-sided buyer-known information is the second-largest effect, not predicted by either paper as a main effect.
\item \textbf{Extended talk backfire} ($-23.3$pp, $p = 0.004$): More communication rounds produce more messages but \emph{worse} outcomes. Likely mechanism: early price anchoring creates commitment that prevents flexible negotiation.
\item \textbf{LLM rapport pooling}: Agents produce generic polite messages (``I'm sure we can find a fair price'') rather than strategic signals. This is a novel behavioral finding about LLM communication patterns.
\end{enumerate}

%% ============================================================
\section{Future Experiment Plan}
%% ============================================================

\paragraph{1. Structured communication.} Replace free-form cheap talk with forced-choice signals (e.g., ``My value is: low / medium / high'') to test whether the pooling result is due to LLM politeness norms or genuine strategic behavior. This isolates the partition-signaling mechanism from CS82.

\paragraph{2. Mechanism study: anchoring vs.\ rapport.} Run a $2 \times 2$ design crossing communication (none vs.\ extended) with anchoring control (price mention allowed vs.\ prohibited in pre-play messages). If prohibiting price mentions eliminates the extended-talk penalty, anchoring is the mechanism.

\paragraph{3. Cross-model replication.} Run the same 12 conditions with Claude Sonnet 4.6 and Gemini 2.5 to test whether the pooling result is model-specific or universal. The pilot (gpt-4o-mini) showed a +17pp cheap talk benefit that vanished in production (gpt-4o), suggesting model capability affects strategic communication.

\paragraph{4. Temperature sensitivity.} Vary agent temperature (0.3, 0.7, 1.0) within a single condition to test whether higher temperature produces more diverse (and potentially more informative) communication strategies.

\paragraph{5. Hypothesis iteration.} Feed these results into \texttt{/hypothesize iterate} with the core finding: ``LLM agents play pooling equilibria universally; extended communication is harmful through anchoring.'' The next hypothesis should test whether structural interventions (forced signals, commitment devices) can break the pooling equilibrium.

%% ============================================================
\section*{Reproducibility Artifacts}
%% ============================================================
\begin{itemize}
\item \texttt{analysis/raw\_sessions.csv} --- 360 rows, 20 columns
\item \texttt{analysis/summary\_by\_condition.csv} --- 12 conditions with Wilson CIs
\item \texttt{analysis/model\_results.csv} --- 19 test rows with effect sizes and $p$-values
\item \texttt{analysis/prior\_alignment.csv} --- 14 prior claims mapped to observed results
\item \texttt{evidence\_trace.csv} --- 12 design choices traced to source papers
\item \texttt{condition\_matrix.csv} --- 12 conditions with factor levels
\item \texttt{results.md} --- Markdown report (6-section format)
\item \texttt{results.tex} --- This document
\item \texttt{results.pdf} --- Compiled PDF
\end{itemize}

\end{document}
